\chapter{Kummer Surfaces and the Orbifold \texorpdfstring{$T^4/\Z_2$}{T4/Z2}}\label{ch:kummer}

The K3 surfaces of \cref{ch:k3} were defined by equations, and their lattice
$3U\oplus2E_8(-1)$ was obtained in \cref{ch:k3lattice} from general theorems. Neither
description lets one \emph{see} the $22$ two-cycles, and neither is a space on which a string
can be quantized exactly. This chapter studies the one K3 surface for which both are possible.
Take a four-torus, identify $x$ with $-x$, and repair the $16$ singular points that result. The
outcome is a K3 surface, the \emph{Kummer surface}\index{Kummer surface}, and every number
attached to K3 can be counted on it by hand: $22=6+16$, $24=8+16$, $20=4+16$,
$(3,19)=(3,3)+(0,16)$, and, for the string, $80=16+64$. The questions we answer are these. Why
are there exactly $16$ singular points and what do they look like? What does ``blowing up''
do, in coordinates, and why does the new sphere have self-intersection $-2$? Why do the $6+16$
obvious classes \emph{not} generate $H^2(X,\Z)$, and what is the rank-$16$ \emph{Kummer
lattice} that they span after saturation? Finally, how does the string know about the $16$
spheres when they have zero size? The answer to the last question, twisted sectors, is the
reason why string theory on an orbifold is consistent although the space is singular. The
torus side of the story uses \cref{ch:narain}; the lattice vocabulary (discriminant, index,
$U$, even lattices) is that of \cref{ch:lattices}.

%=====================================================================================
\section{The orbifold \texorpdfstring{$T^4/\Z_2$}{T4/Z2}}\label{sec:kummer:orbifold}

\subsection{Sixteen fixed points}

Let $\Lambda\subset\R^4$ be a lattice of rank four and $A=\R^4/\Lambda$ the four-torus. The map
$\iota\colon x\mapsto -x$ preserves $\Lambda$ and therefore acts on $A$; since
$\iota^2=1$ it generates a group $\Z_2$. We write $S_0=A/\iota$ for the quotient, in which the
points $x$ and $-x$ are identified, and follow the notation of \source{Asp §2.6,
ll.~1064--1072}. A point of $A$ is fixed by $\iota$ when $x\equiv-x$ modulo $\Lambda$, that is
when $2x\in\Lambda$:
\begin{equation}\label{eq:kummer:fixed}
  \mathrm{Fix}(\iota)=\tfrac12\Lambda/\Lambda\;\cong\;(\Z/2\Z)^4 ,\qquad
  |\mathrm{Fix}(\iota)|=2^4=16 .
\end{equation}
If $\lambda_1,\dots,\lambda_4$ is a basis of $\Lambda$, the fixed points are
$p_\epsilon=\frac12(\epsilon_1\lambda_1+\dots+\epsilon_4\lambda_4)$ with
$\epsilon=(\epsilon_1,\dots,\epsilon_4)\in\mathbb{F}_2^4$, where $\mathbb{F}_2=\{0,1\}$ is the
field with two elements. These are the points of order two of the group $A$
\source{Huy ch.~1, Example~1.3(iii), ll.~218--222}. The labelling by the vector space
$\mathbb{F}_2^4$\index{two-torsion points} looks like bookkeeping now; in
\cref{sec:kummer:lattice} the affine geometry of $\mathbb{F}_2^4$ will decide which
combinations of cycles are integral.

\begin{example}[The square torus]\label{ex:kummer:square}
Take $A=\C^2/(\Z+\ii\Z)^2$ with coordinates $(z_1,z_2)$, as in \source{Asp eq.~(49),
ll.~1066--1072}. In each factor the fixed points of $z\mapsto-z$ are
$0,\ \frac12,\ \frac{\ii}2,\ \frac12+\frac{\ii}2$, so the $16$ fixed points are the pairs
$(z_1,z_2)$ with both entries in this list. One factor alone gives the ``pillow''
$T^2/\Z_2$ of \cref{fig:kummer:pillow}: fold the fundamental square along its middle line and
sew the edges; the result is a sphere with four corners.
\end{example}

\begin{figure}[t]
\centering
\begin{tikzpicture}[scale=0.82]
  % pillow: fundamental domain
  \fill[blue!8] (0,0) rectangle (1.5,3);
  \draw[thick] (0,0) rectangle (3,3);
  \draw[dashed] (1.5,0)--(1.5,3);
  \foreach \p in {(0,0),(1.5,0),(0,1.5),(1.5,1.5)} \fill[tierC] \p circle (2.2pt);
  \foreach \p in {(3,0),(0,3),(3,3),(1.5,3),(3,1.5)} \draw[tierC] \p circle (2.2pt);
  \node[below] at (0,0) {\small $0$};
  \node[below] at (1.5,0) {\small $\frac12$};
  \node[left] at (0,1.5) {\small $\frac{\ii}2$};
  \node[below] at (3,0) {\small $1$};
  \node[left] at (0,3) {\small $\ii$};
  \node at (1.5,-0.95) {\small (a) $T^2/\Z_2$};
  % 4x4 grid = F_2^4
  \begin{scope}[xshift=5.2cm]
    \fill[green!12,rounded corners=3pt] (-0.3,-0.3) rectangle (1.3,3.3);
    \draw[tierA,thick,rounded corners=3pt] (-0.22,-0.22) rectangle (3.22,0.22);
    \foreach \x in {0,1,2,3} \foreach \y in {0,1,2,3} \fill (\x,\y) circle (2pt);
    \node[below] at (0,-0.3) {\scriptsize $00$};
    \node[below] at (1,-0.3) {\scriptsize $01$};
    \node[below] at (2,-0.3) {\scriptsize $10$};
    \node[below] at (3,-0.3) {\scriptsize $11$};
    \node[left] at (-0.3,0) {\scriptsize $00$};
    \node[left] at (-0.3,1) {\scriptsize $01$};
    \node[left] at (-0.3,2) {\scriptsize $10$};
    \node[left] at (-0.3,3) {\scriptsize $11$};
    \node at (1.5,-0.95) {\small (b) $\mathbb{F}_2^4$: $(\epsilon_1\epsilon_2\,|\,\epsilon_3\epsilon_4)$};
  \end{scope}
\end{tikzpicture}
\caption{(a) A fundamental domain of $z\mapsto-z$ on the square torus is the shaded half; the
four filled dots are the fixed points (open dots are their translates). (b) The $16$ fixed
points of $T^4/\Z_2$ as the points of $\mathbb{F}_2^4$; horizontal position is
$(\epsilon_1\epsilon_2)$, vertical position is $(\epsilon_3\epsilon_4)$. The framed row is an
affine plane (four points: the fixed points on one invariant two-torus); the shaded block
$\epsilon_1=0$ is an affine hyperplane (eight points).}\label{fig:kummer:pillow}
\end{figure}

\subsection{What the quotient looks like near a fixed point}

Away from the $16$ points the action of $\Z_2$ is free and $S_0$ is a smooth flat manifold.
Near a fixed point $p$ put $y=x-p$; then $\iota$ acts as $y\mapsto-y$, and a neighbourhood of
the image of $p$ in $S_0$ is a neighbourhood of the origin in
\begin{equation}\label{eq:kummer:local}
  \R^4/\{\pm1\}=\C^2/\Z_2,\qquad (\xi,\eta)\sim(-\xi,-\eta).
\end{equation}
This is a cone: the sphere $|y|=r$ is mapped to $S^3/\{\pm1\}=\mathbb{RP}^3$, so
$\C^2/\Z_2$ is the cone over $\mathbb{RP}^3$ with its vertex at the fixed point. It is not a
manifold there, because a small punctured neighbourhood of a point in a four-manifold is
simply connected, while $\pi_1(\mathbb{RP}^3)=\Z_2$. The space $S_0$ is an
\emph{orbifold}\index{orbifold}: it is covered by patches of the form $\R^n/G_i$ with $G_i$
finite \source{Asp §2.6, ll.~1042--1049}. Such a singular point, locally $\C^2/\Z_2$, is
called an $A_1$ singularity\index{A1 singularity@$A_1$ singularity} or ordinary double point;
the name will be justified in \cref{sec:kummer:local}.

\subsection{Why \texorpdfstring{$S_0$}{S0} deserves to be called a K3 surface}

A K3 surface was defined in \cref{ch:k3} by a trivial canonical bundle and the absence of
holomorphic one-forms, and was shown to be simply connected. We check the three properties
for $S_0$. Choose a complex structure on $\R^4$ with coordinates $(z_1,z_2)$, so that $A$ is
a complex torus.
\begin{enumerate}[leftmargin=*,itemsep=2pt]
\item The holomorphic two-form $\dd z_1\wedge\dd z_2$ is invariant, since both factors change
sign. It descends to $S_0$ and has no zeros. The Kähler form
$\frac{\ii}2(\dd z_1\wedge\dd\bar z_1+\dd z_2\wedge\dd\bar z_2)$ is invariant for the same
reason \source{Asp ll.~1073--1077}.
\item The one-forms $\dd z_1,\dd z_2$ are odd and do not descend: $S_0$ has no holomorphic
one-forms, although $A$ has two.
\item $S_0$ is simply connected \source{Asp ll.~1074--1076}. Indeed, let $\gamma$ be the
closed loop $s\mapsto s\lambda_1$, $s\in[0,1]$, which generates one factor of
$\pi_1(A)=\Z^4$. Its second half, from the fixed point $\frac12\lambda_1$ to
$\lambda_1\equiv0$, is the image under $\iota$ of the first half run backwards. In $S_0$ the
loop is therefore of the form $\delta\cdot\delta^{-1}$ and can be contracted. The new closed
loops of $S_0$, images of paths from $y$ to $-y$, are of the same form once the path is slid
through a fixed point.
\end{enumerate}

\begin{physicsbox}[Holonomy $\Z_2\subset SU(2)$: a flat space that breaks half the supersymmetry]
The torus is flat, so parallel transport around a contractible loop is trivial. In $S_0$ there
are new closed loops: the image of a path in $A$ from $y$ to $-y$. A vector transported along
it comes back multiplied by $-1$, because the identification $y\sim-y$ acts on tangent
vectors too. The holonomy group of $S_0$ is thus $\Z_2=\{\pm\mathbf 1_2\}$, and
$-\mathbf 1_2=\diag(-1,-1)$ has determinant $+1$: it lies in $SU(2)$
\source{Asp ll.~1331--1335}. A spinor on a four-dimensional space transforms under
$SU(2)_+\times SU(2)_-$; if the holonomy lies in one factor, the spinors of the other
chirality are untouched (\cref{ch:k3}). Hence $S_0$ preserves exactly the same half of the
supersymmetry as a smooth K3 surface with its full $SU(2)$ holonomy, while being flat
everywhere except at $16$ points. All the curvature of K3, in particular the whole of
$\chi=24$, has been pushed into these points.
\end{physicsbox}

\begin{pitfallbox}[$T^2/\Z_2$ is not a ``two-dimensional K3'']
For one complex dimension, $z\mapsto-z$ has determinant $-1$ on $\dd z$: the holomorphic
one-form is projected out, and the holonomy $-1\in U(1)$ is not in $SU(1)=\{1\}$. The pillow
of \cref{fig:kummer:pillow}(a) is a sphere with four conical points of deficit angle $\pi$;
the total deficit $4\pi=2\pi\chi(S^2)$ is the Gauss--Bonnet theorem. It is positively
curved and preserves no supersymmetry by itself. The rule is that
$\C^n/G$ is a local model of a Calabi--Yau orbifold only when $G\subset SU(n)$; for $n=2$
this is what makes $\dd z_1\wedge\dd z_2$ invariant \source{Asp ll.~1331--1333}.
\end{pitfallbox}

%=====================================================================================
\section{The local model: resolving \texorpdfstring{$\C^2/\Z_2$}{C2/Z2}}\label{sec:kummer:local}

\subsection{The singularity as a quadric cone}

Functions on $\C^2/\Z_2$ are functions of $(\xi,\eta)$ that are even under
$(\xi,\eta)\mapsto(-\xi,-\eta)$. Every even polynomial is a polynomial in the three quadratic
monomials
\begin{equation}\label{eq:kummer:invariants}
  x_0=\xi^2,\qquad x_1=\eta^2,\qquad x_2=\xi\eta,\qquad\text{which satisfy}\qquad
  f:=x_0x_1-x_2^2=0 .
\end{equation}
Conversely a point of the surface $\{f=0\}\subset\C^3$ determines $(\xi,\eta)$ up to the
common sign, so $\C^2/\Z_2$ \emph{is} the hypersurface $f=0$ \source{Asp eq.~(50),
ll.~1082--1091}. The gradient $(\partial_0f,\partial_1f,\partial_2f)=(x_1,x_0,-2x_2)$ vanishes
only at the origin, which is the single singular point. With $x_0=u+\ii v$, $x_1=u-\ii v$,
$x_2=\ii w$ the equation becomes $u^2+v^2+w^2=0$, the case $n=2$ of the $A_{n-1}$ singularity $x^2+y^2+z^n=0$ in the $A$--$D$--$E$ list\index{ADE singularities@$A$--$D$--$E$ singularities} of
surface singularities \source{Asp ll.~1226--1230 and eq.~(58), ll.~1523--1527}.

\subsection{Blowing up in coordinates}\index{blow-up}

The idea of a blow-up is to remember not only \emph{that} a path reaches the singular point
but \emph{from which direction}. In $\C^3$ a direction is a point $[s_0:s_1:s_2]\in
\mathbb{P}^2$, and the blow-up of $\C^3$ at the origin is the set of pairs (point, line
through the origin containing it), $\{x_is_j=x_js_i\}\subset\C^3\times\mathbb{P}^2$
\source{Asp eq.~(51), ll.~1094--1106}. A path in the cone $f=0$ approaching the origin along
$(x_0t,x_1t,x_2t)$ lands on a direction with $s_0s_1-s_2^2=0$. The closure of the cone minus
its vertex inside the blow-up, its \emph{proper transform} $\widetilde{\C^2/\Z_2}$, therefore
replaces the vertex by the conic
\begin{equation}\label{eq:kummer:E}
  E=\{s_0s_1=s_2^2\}\subset\mathbb{P}^2,\qquad
  [\xi:\eta]\longmapsto[\xi^2:\eta^2:\xi\eta],
\end{equation}
which the displayed parametrization identifies with $\mathbb{P}^1$
\source{Asp ll.~1123--1131}. In words: $E$ is the set of complex lines $\ell\subset\C^2$
through the origin (the points $\pm(\xi,\eta)$ lie on the same line), and the resolved space
is
\begin{equation}\label{eq:kummer:total}
  \widetilde{\C^2/\Z_2}=\{(\ell,q)\;:\;\ell\in\mathbb{P}^1,\ q\in\ell/\{\pm1\}\} .
\end{equation}
It is fibred over $E=\mathbb{P}^1$ with fibre $\ell/\{\pm1\}\cong\C$ (the quotient of a
complex line by $\pm1$ is again a complex line, with coordinate the square). We now make this
explicit. Cover $\mathbb{P}^1$ by the charts $\xi\neq0$ and $\eta\neq0$ and use
\begin{equation}\label{eq:kummer:charts}
  \text{chart 1:}\quad t=\frac{\eta}{\xi}=\frac{x_2}{x_0},\quad a=\xi^2=x_0;
  \qquad\qquad
  \text{chart 2:}\quad t'=\frac{\xi}{\eta}=\frac{x_2}{x_1},\quad a'=\eta^2=x_1 .
\end{equation}
In chart 1 the cone is parametrized by $(x_0,x_1,x_2)=(a,\,at^2,\,at)$, which indeed solves
$x_0x_1=x_2^2$, and $(a,t)\in\C^2$ are unconstrained coordinates: the space is smooth. The
exceptional curve is $E=\{a=0\}$, with $t$ as its coordinate. On the overlap
\begin{equation}\label{eq:kummer:transition}
  t'=\frac1t,\qquad a'=t^2\,a .
\end{equation}
Compare with the tautological line bundle $\mathcal{O}(-1)$ over $\mathbb{P}^1$, whose fibre
over $\ell$ is the line $\ell$ itself: its fibre coordinates are $\xi$ and $\eta=t\,\xi$, so
the transition function is $t$. The transition function of our bundle is $t^2$, hence
\begin{equation}\label{eq:kummer:Ominus2}
  \widetilde{\C^2/\Z_2}\;=\;\text{total space of }\mathcal{O}_{\mathbb{P}^1}(-2),
\end{equation}
with $E$ as the zero section \source{Asp ll.~1134--1137}.

\begin{proposition}[The exceptional curve is a $(-2)$-curve]\label{prop:kummer:minus2}
$E\cdot E=-2$.\index{(-2)-curve@$(-2)$-curve}
\end{proposition}
\begin{proof}[Two derivations]
(i) The self-intersection of a curve is the degree of its normal bundle, computed as the
number of zeros minus the number of poles of any meromorphic section \source{Asp
ll.~1139--1173}. Take the section $a=1$ in chart 1. In chart 2 it reads
$a'=t^2=1/t'^2$: no zeros and a double pole at $t'=0$. The degree is $0-2=-2$.
(ii) For a smooth curve $C$ of genus $g$ in a surface with trivial canonical bundle, the
adjunction formula gives $C\cdot C=2g-2$ \source{Asp eq.~(53), ll.~1178--1182}. For
$g=0$ this is $-2$. Route (ii) assumes the canonical bundle of the resolved space to be
trivial; this is \cref{prop:kummer:crepant}.
\end{proof}

A negative self-intersection means that $E$ is rigid: it cannot be moved to a nearby
holomorphic curve, since two distinct holomorphic curves intersect non-negatively
\source{Asp ll.~1173--1177}.

\begin{proposition}[The resolution is crepant]\label{prop:kummer:crepant}
The two-form $\dd\xi\wedge\dd\eta$ on $\C^2/\Z_2$ extends to a holomorphic two-form without
zeros on $\widetilde{\C^2/\Z_2}$. Explicitly
$\dd\xi\wedge\dd\eta=\tfrac12\,\dd a\wedge\dd t=-\tfrac12\,\dd a'\wedge\dd t'$.
\index{crepant resolution}
\end{proposition}
\begin{proof}
In chart 1, $\xi=a^{1/2}$ and $\eta=t\,a^{1/2}$, so
$\dd\xi=\frac12a^{-1/2}\dd a$ and $\dd\eta=a^{1/2}\dd t+\frac12ta^{-1/2}\dd a$. The wedge
product is $\frac12a^{-1/2}a^{1/2}\,\dd a\wedge\dd t=\frac12\dd a\wedge\dd t$: the square roots
cancel, and nothing special happens at $a=0$. From \eqref{eq:kummer:transition},
$\dd a'\wedge\dd t'=(t^2\dd a+2ta\,\dd t)\wedge(-t^{-2}\dd t)=-\dd a\wedge\dd t$. (Both Jacobians
were confirmed by a symbolic check.)
\end{proof}

\begin{pitfallbox}[Blowing up a smooth point is different]
If one blows up the origin of $\C^2$ itself, without the quotient, the coordinates are
$(\xi,t)$ with $\eta=t\xi$, the exceptional curve has normal bundle $\mathcal{O}(-1)$ and
self-intersection $-1$, and $\dd\xi\wedge\dd\eta=\xi\,\dd\xi\wedge\dd t$ \emph{vanishes} along
$E=\{\xi=0\}$. The canonical bundle changes, $\omega_{\tilde A}\cong\mathcal{O}(\sum E_i)$ for
the blow-up $\tilde A$ of the torus in $16$ points \source{Huy ch.~1, ll.~244--246}. Blowing
up a smooth point of a K3 surface destroys the K3 property; blowing up an $A_1$ point
preserves it.
\end{pitfallbox}

\subsection{The global construction, in two orders}

Performing the local surgery at all $16$ points of $S_0$ gives a smooth compact complex
surface $X=\operatorname{Km}(A)$ with a nowhere vanishing holomorphic two-form
(\cref{prop:kummer:crepant} glued to $\dd z_1\wedge\dd z_2$), with no holomorphic one-forms
and simply connected: a K3 surface \TL{} \source{Asp ll.~1188--1194; Huy ch.~1,
Example~1.3(iii), ll.~218--250}. One can also resolve in the opposite order
\source{Huy ll.~225--240 and §3.2.5, ll.~2015--2031}: first blow up the $16$ fixed points of
the smooth torus, obtaining $\sigma\colon\tilde A\to A$ with sixteen $(-1)$-curves $E_i$; the
involution lifts to $\tilde\iota$ on $\tilde A$, acting in the coordinates of the pitfall box
as $(\xi,t)\mapsto(-\xi,t)$; its fixed locus is the union of the curves $E_i=\{\xi=0\}$, and
the quotient has the smooth coordinates $(a,t)=(\xi^2,t)$ of \eqref{eq:kummer:charts}:
\begin{equation}\label{eq:kummer:square}
\begin{tikzcd}[column sep=large]
  \tilde A \arrow[r,"\sigma"] \arrow[d,"\pi"'] & A \arrow[d] \\
  X \arrow[r] & A/\iota=S_0
\end{tikzcd}
\qquad\qquad
\begin{aligned}
  &\pi\colon\tilde A\to X \text{ is a double cover branched along } \textstyle\bigcup E_i,\\
  &\bar E_i:=\pi(E_i)\subset X,\qquad \pi^*\mathcal{O}(\bar E_i)\cong\mathcal{O}(2E_i).
\end{aligned}
\end{equation}
The second order is the one used for computations in cohomology, because everything is
pulled back from or pushed forward to the smooth surface $\tilde A$.

\begin{figure}[t]
\centering
\begin{tikzpicture}[scale=0.72]
  % singular cone
  \draw[thick] (-1.5,2)--(0,0)--(1.5,2);
  \draw[thick] (-1.5,-2)--(0,0)--(1.5,-2);
  \draw (0,2) ellipse (1.5 and 0.3);
  \draw (1.5,-2) arc (0:-180:1.5 and 0.3);
  \draw[dashed] (1.5,-2) arc (0:180:1.5 and 0.3);
  \fill[tierC] (0,0) circle (2.5pt);
  \node[right] at (0.25,0) {\small $A_1$ point};
  \node at (0,-2.9) {\small $\C^2/\Z_2$: \ $x_0x_1=x_2^2$};
  % arrow
  \draw[-{Stealth},thick] (5.0,0)--(2.9,0) node[midway,above] {\small blow down};
  \node at (3.95,-0.4) {\small $J\cdot E\to0$};
  % resolved
  \begin{scope}[xshift=7.5cm]
    \draw[thick] (-1.5,2) .. controls (-0.5,0.7) and (-0.5,-0.7) .. (-1.5,-2);
    \draw[thick] (1.5,2) .. controls (0.5,0.7) and (0.5,-0.7) .. (1.5,-2);
    \draw (0,2) ellipse (1.5 and 0.3);
    \draw (1.5,-2) arc (0:-180:1.5 and 0.3);
    \draw[dashed] (1.5,-2) arc (0:180:1.5 and 0.3);
    \draw[tierA,very thick] (0.75,0) arc (0:-180:0.75 and 0.18);
    \draw[tierA,very thick,dashed] (0.75,0) arc (0:180:0.75 and 0.18);
    \node[right] at (0.85,0) {\small $E\cong\mathbb{P}^1$,\ $E^2=-2$};
    \node at (0,-2.9) {\small total space of $\mathcal{O}_{\mathbb{P}^1}(-2)$};
  \end{scope}
\end{tikzpicture}
\caption{The $A_1$ singularity and its resolution (real two-dimensional cartoon of a complex
surface). The vertex of the cone is replaced by a sphere $E$ whose area is $J\cdot E$, the
integral of the Kähler form. Letting $J\cdot E\to0$ with the complex structure fixed recovers
the orbifold.}\label{fig:kummer:cone}
\end{figure}

\begin{physicsbox}[The size of $E$ is a modulus]
Nothing in the algebraic construction fixes the \emph{size} of the new sphere. With a Kähler
form $J$ on $X$ the area of $\bar E_i$ is $J\cdot\bar E_i=\int_{\bar E_i}J$. Keeping the
complex structure fixed, one may move $J$ until it becomes orthogonal to $\bar E_i$; the
sphere then has zero area and the surface is the orbifold again \source{Asp
ll.~1195--1202}. The orbifold $S_0$ is thus not a different space from the Kummer surface
but a boundary point of its family of metrics: in the language of \cref{ch:k3lattice}, a
Ricci-flat metric is a positive three-plane $\Sigma\subset H^2(X,\R)$, and the metric is an
orbifold metric exactly when $\Sigma$ is orthogonal to a root, a class $\alpha$ with
$\alpha^2=-2$ \TL{} \source{Asp Theorem~4, ll.~1208--1216}. At $S_0$ the plane $\Sigma$ is
orthogonal to all sixteen $\bar E_i$ at once. Near each resolved point the Ricci-flat metric
is approximated by an explicitly known asymptotically flat metric on
$\mathcal{O}_{\mathbb{P}^1}(-2)$ \source{Asp ll.~1503--1507}. Because $\Sigma$ is spanned by
$J$ and the real and imaginary parts of the holomorphic two-form $\Omega$, one can also make
$\bar E_i$ disappear by rotating $\Omega$ instead of $J$: the singularity is then
\emph{deformed}, $x_0x_1-x_2^2=\varepsilon$, rather than blown up, and for K3 the two
procedures are related by a rotation of $\Sigma$ \source{Asp ll.~1513--1530}. Each
$A_1$ point therefore carries three real metric parameters, the components of $\Sigma$ along
$\bar E_i$.
\end{physicsbox}

%=====================================================================================
\section{Counting: \texorpdfstring{$b_2=6+16$ and $\chi=24$}{b2 = 6+16 and chi = 24}}
\label{sec:kummer:count}

\subsection{The cohomology of the orbifold}

With rational coefficients the cohomology of a quotient by a finite group is the invariant
part of the cohomology of the covering space (average a closed form over the group). For the
torus, $H^k(A,\Q)=\Lambda^kH^1(A,\Q)$ is spanned by the products
$\dd x^{m_1}\wedge\dots\wedge\dd x^{m_k}$, and $\iota^*\dd x^m=-\dd x^m$. Therefore
$\iota^*=(-1)^k$ on $H^k$ \source{Huy §3.2.5, ll.~2037--2038}, and
\begin{equation}\label{eq:kummer:bettiS0}
  b_k(S_0)=\tbinom4k\ \text{for $k$ even},\quad 0\ \text{for $k$ odd}:
  \qquad (b_0,\dots,b_4)=(1,0,6,0,1),\qquad \chi(S_0)=8 .
\end{equation}
The six surviving two-forms $\dd x^m\wedge\dd x^n$ are dual to the six coordinate two-tori
$T_{mn}\subset A$. The Hodge decomposition of the invariant forms is just as simple. On the
torus $h^{p,q}(A)=\binom2p\binom2q$, spanned by
$\dd z_I\wedge\dd\bar z_J$, and the form has sign $(-1)^{p+q}$ under $\iota$:
\begin{equation}\label{eq:kummer:hodgeS0}
  h^{2,0}=h^{0,2}=1\ (\dd z_1\wedge\dd z_2\text{ and its conjugate}),\qquad
  h^{1,1}=4\ (\dd z_i\wedge\dd\bar z_j),\qquad h^{1,0}=h^{0,1}=0 .
\end{equation}

\subsection{Adding the sixteen spheres}

Let $N_i$ be a small conical neighbourhood of the $i$-th singular point and $\tilde N_i$ the
neighbourhood of $\bar E_i$ that replaces it. $N_i$ is contractible; $\tilde N_i$ is a disc
bundle over $S^2$ and contracts onto $\bar E_i$. Their common boundary is $\mathbb{RP}^3$,
which has the rational cohomology of $S^3$. Writing $X=(X\setminus\bigcup\tilde
N_i)\cup\bigcup\tilde N_i$, the Mayer--Vietoris sequence shows that each replacement adds
exactly one class to $H^2(\,\cdot\,,\Q)$, the class of $\bar E_i$, and changes nothing in
degrees $1$ and $3$. Hence
\begin{equation}\label{eq:kummer:b2}
  b_2(X)=6+16=22,\qquad b_1(X)=b_3(X)=0,\qquad \chi(X)=1+22+1=24 .
\end{equation}
Since the curves $\bar E_i$ are holomorphic, their classes are of type $(1,1)$, and
\eqref{eq:kummer:hodgeS0} becomes $h^{1,1}(X)=4+16=20$, $h^{2,0}(X)=1$: the Hodge diamond of
\cref{ch:k3}. The same grouping, $16$ classes from the singular points, the $(2,0)$ and
$(0,2)$ forms and $4$ more $(1,1)$ forms from the torus, is the one used in the physics
literature \source{DRS §2, ll.~513--519}. For the signature, the wedge product on the
invariant classes $\Lambda^2\R^4$ has signature $(3,3)$ (self-dual and anti-self-dual forms),
and the sixteen classes $\bar E_i$ are mutually orthogonal with square $-2$:
\begin{equation}\label{eq:kummer:signature}
  (3,3)+(0,16)=(3,19),
\end{equation}
in agreement with the signature found from the Hirzebruch theorem in \cref{ch:k3}.

\begin{example}[The number $24$ in three ways]\label{ex:kummer:24}
\emph{(a) Betti numbers:} \eqref{eq:kummer:b2}. \ \emph{(b) Cut and paste:} the Euler
characteristic is additive under decompositions and multiplicative under unbranched
covers. Remove the $16$ fixed points from $A$: $\chi=0-16$. The quotient of the rest is free:
$\chi=-8$. Glue in $16$ spheres: $\chi(X)=-8+16\cdot2=24$. Gluing in $16$ points instead gives
$\chi(S_0)=-8+16=8$, as in \eqref{eq:kummer:bettiS0}.
\ \emph{(c) Elliptic fibration:} let $A=E_1\times E_2$ be a product of elliptic curves. The
projection to $E_1/\pm\cong\mathbb{P}^1$ makes $X$ an elliptic surface whose smooth fibres
are copies of $E_2$ ($\chi=0$). Over each of the four branch points the fibre of
$S_0$ is $E_2/\pm\cong\mathbb{P}^1$ and passes through four singular points; in $X$ it becomes
a configuration of five spheres, the central one meeting the four $\bar E_i$ once each: a
fibre of Kodaira type $I_0^*$\index{elliptic fibration!$I_0^*$ fibre} \source{Huy ch.~11, Example~1.11(i), ll.~9500--9507}. Five
spheres glued at four points have $\chi=5\cdot2-4=6$, and $\chi(X)=4\cdot6=24$; equivalently,
the discriminant of an elliptic K3 surface has degree $24$ and vanishes to order $6$ at an
$I_0^*$ fibre \source{Huy §11.2.4, ll.~9725--9735}.
All three counts agree with $c_2(X)=24$ from Noether's formula in
\cref{ch:k3}.
\end{example}

%=====================================================================================
\section{The intersection form and the Kummer lattice}\label{sec:kummer:lattice}

The count $22=6+16$ is a statement about vector spaces. The integral statement is subtler
and is the mathematical core of this chapter: the $6+16$ obvious classes span a sublattice of
index $2^{11}$ in $H^2(X,\Z)$.

\subsection{The torus classes are rescaled by two}

Let $\alpha,\alpha'\in H^2(\tilde A,\Z)$ be $\tilde\iota$-invariant classes and
$\pi_*\alpha,\pi_*\alpha'$ their push-forwards to $X$. Since $\pi$ has degree two and
$\alpha$ is invariant, $\pi^*\pi_*\alpha=2\alpha$. On the other hand
$(\pi^*\beta\cdot\pi^*\beta')_{\tilde A}=2\,(\beta\cdot\beta')_X$ for all classes on $X$,
again because $\pi$ has degree two. Combining the two, $4(\alpha\cdot\alpha')_{\tilde
A}=2(\pi_*\alpha\cdot\pi_*\alpha')_X$, that is
\begin{equation}\label{eq:kummer:factor2}
  (\pi_*\alpha\cdot\pi_*\alpha')_X=2\,(\alpha\cdot\alpha')_{\tilde A}
  \qquad\text{\source{Huy eq.~(2.5), ll.~2037--2044}} .
\end{equation}
Two consequences follow. First, $\pi_*[E_i]=[\bar E_i]$ and $E_i^2=-1$ give
$\bar E_i^2=-2$: a third derivation of \cref{prop:kummer:minus2}. Second, $H^2(A,\Z)$ with
the basis $\dd x^m\wedge\dd x^n$ is $U^{\oplus3}$, the three hyperbolic pairs being
$(T_{12},T_{34})$, $(T_{13},T_{42})$, $(T_{14},T_{23})$ with $T_{12}\cdot T_{34}=1$ and so
on. Its image in $H^2(X,\Z)$ carries twice this form:
\begin{equation}\label{eq:kummer:U2}
  \pi_*H^2(A,\Z)\;\cong\;U(2)^{\oplus3},\qquad
  \operatorname{disc}U(2)^{\oplus3}=\det\begin{pmatrix}0&2\\2&0\end{pmatrix}^3=-2^6 .
\end{equation}
The factor two has a transparent meaning. A two-torus $T_{12}$ in generic position and its
mirror image $\iota(T_{12})$ are different in $A$ and become one torus in $X$. The images of
$T_{12}$ and $T_{34}$ meet where $(T_{12}\cup\iota T_{12})\cap(T_{34}\cup\iota T_{34})$, four
points of $A$, is mapped to: two points of $X$.

\subsection{Half-integral classes from invariant tori}

\begin{lemma}\label{lem:kummer:C}
Let $T\subset A$ be a two-torus through a fixed point, invariant under $\iota$, for instance
$T_{12}$ placed at $(x^3,x^4)=(0,0)$. It contains four fixed points, forming an affine plane
$P\subset\mathbb{F}_2^4$. The proper transform in $X$ of the image of $T$ is a sphere $C$
with
\begin{equation}\label{eq:kummer:C}
  [C]=\tfrac12\Bigl(\pi_*[T]-\sum_{i\in P}[\bar E_i]\Bigr)\in H^2(X,\Z),\qquad
  C^2=-2,\qquad C\cdot\bar E_i=1\ \ (i\in P).
\end{equation}
\end{lemma}
\begin{proof}
The statement is topological, so we may choose the complex structure on $A$ for which $T$ is
a holomorphic curve ($z_1=x^1+\ii x^2$ for $T_{12}$). $T/\iota$ is the pillow of \cref{fig:kummer:pillow}(a), a sphere through four singular
points. In $\tilde A$ the proper transform $\tilde T$ of $T$ has class
$\sigma^*[T]-\sum_{i\in P}[E_i]$, since $T$ passes once through each blown-up point. The map
$\tilde T\to C$ is two-to-one, so $\pi_*[\tilde T]=2[C]$; this is the first formula. Then, by
$\pi_*[T]^2=2\,T^2=0$ and the orthogonality of $\pi_*[T]$ to the $\bar E_i$,
$C^2=\frac14\bigl(0+4\cdot(-2)\bigr)=-2$ and
$C\cdot\bar E_i=\frac12\cdot(-1)\cdot(-2)=1$. The value $C^2=-2$ is what \emph{every} smooth
sphere in a K3 surface must have, which checks the formula.
\end{proof}

So $\frac12\bigl(\pi_*[T]-\sum_P\bar E_i\bigr)$ is an integral class although neither half is:
the sublattice $\pi_*H^2(A,\Z)\oplus\bigoplus\Z\bar E_i$ is not all of $H^2(X,\Z)$. The class
$F=2C+\sum_{i\in P}\bar E_i=\pi_*[T]$ is the $I_0^*$ fibre of \cref{ex:kummer:24}(c); its five
components form the affine Dynkin diagram $\tilde D_4$ of \cref{fig:kummer:D4}, and
$F^2=0$ as it must be for a fibre.

\begin{figure}[t]
\centering
\begin{tikzpicture}[scale=0.8,every node/.style={font=\small}]
  \node[circle,draw,thick,inner sep=2pt] (c) at (0,0) {$2$};
  \foreach \ang/\lab in {45/{\bar E_{00}},135/{\bar E_{01}},225/{\bar E_{10}},315/{\bar E_{11}}}
    {\node[circle,draw,thick,inner sep=2pt] (e\ang) at (\ang:1.5) {$1$};
     \draw[thick] (c)--(e\ang);
     \node at (\ang:2.25) {$\lab$};}
  \node at (0,-0.55) {$C$};
  \begin{scope}[xshift=6.6cm]
   \draw[thick] (-2,0)--(2,0); \node[right] at (2,0) {$C$};
   \foreach \x/\lab in {-1.5/{00},-0.5/{01},0.5/{10},1.5/{11}}
     {\draw[thick,tierA] (\x,-1)--(\x,1); \node[above] at (\x,1) {$\bar E_{\lab}$};}
  \end{scope}
\end{tikzpicture}
\caption{The image of an invariant two-torus. Left: dual graph $\tilde D_4$; nodes are
$(-2)$-spheres, edges are transverse intersection points, the numbers are the multiplicities
in the fibre class $F=2C+\sum\bar E_i$. Right: the curves themselves. The labels are the
$(\epsilon_1\epsilon_2)$ of \cref{fig:kummer:pillow}(b).}\label{fig:kummer:D4}
\end{figure}

\subsection{From planes to hyperplanes: the Kummer lattice}

Take two \emph{parallel} invariant tori $T,T'$ (the same $T_{12}$ placed at two different
fixed values of $(x^3,x^4)$), with planes $P,P'$. They are homologous, so subtracting the
two classes \eqref{eq:kummer:C} eliminates the torus part:
\begin{equation}\label{eq:kummer:halfW}
  [C']-[C]=\tfrac12\sum_{i\in P}\bar E_i-\tfrac12\sum_{i\in P'}\bar E_i
  \;\equiv\;\tfrac12\sum_{i\in W}\bar E_i\pmod{\textstyle\bigoplus\Z\bar E_i},
  \qquad W=P\cup P' .
\end{equation}
The union of two parallel affine planes in $\mathbb{F}_2^4$ is an affine hyperplane, a set of
$8$ points defined by one equation $a\cdot\epsilon=c$ (in \cref{fig:kummer:pillow}(b): two
rows). Every affine hyperplane arises in this way: choose a basis of $\Lambda$ adapted to
$W$ and take two parallel invariant tori spanned by two of the basis vectors.
Hence, for each of the $30$ affine hyperplanes $W\subset\mathbb{F}_2^4$, the class
$\frac12\sum_{i\in W}\bar E_i$ lies in $H^2(X,\Z)$; adding two complementary hyperplanes, so
does $\frac12\sum_{\text{all }16}\bar E_i$, in agreement with the square root $L$ of
$\mathcal{O}(\sum\bar E_i)$ that defines the double cover $\pi$ \source{Huy ch.~1,
ll.~249--251}.

\begin{definition}[Kummer lattice]\label{def:kummer:K}\index{Kummer lattice}
The Kummer lattice $K$ is the saturation of $\bigoplus_i\Z\bar E_i$ in $H^2(X,\Z)$, i.e.\ the
smallest primitive sublattice containing all $\bar E_i$, equivalently
$K=(\bigoplus\Z\bar E_i)^{\perp\perp}$ \source{Huy Def.~14.3.13, ll.~13469--13476}.
(A sublattice $M\subseteq\Lambda$ is \emph{primitive}\index{primitive sublattice} if
$\Lambda/M$ has no torsion: $nv\in M$ with $v\in\Lambda$ implies $v\in M$.)
\end{definition}

\begin{theorem}[Structure of the Kummer lattice; \TL]\label{thm:kummer:K}
$K$ is generated by the $\bar E_i$ together with the $\frac12\sum_{i\in W}\bar E_i$, $W$ an
affine hyperplane of $\mathbb{F}_2^4$ \source{Huy §14.3.3, ll.~13482--13487}. Moreover
\source{Huy Prop.~14.3.14, ll.~13497--13503}:
\begin{enumerate}[label=(\roman*),itemsep=1pt]
\item $K^\perp=\pi_*H^2(A,\Z)\cong U(2)^{\oplus3}$;
\item the index of $\bigoplus\Z\bar E_i\cong\langle-2\rangle^{\oplus16}$ in $K$ is $2^5$;
\item $K$ is even, negative definite, of rank $16$ and $\operatorname{disc}K=2^6$;
\item the discriminant group is $A_K\cong(\Z/2\Z)^6$.
\end{enumerate}
$K$ does not depend on the torus $A$, and its roots (vectors of square $-2$) are exactly the
$\pm\bar E_i$ \source{Huy ll.~13494--13496}.
\end{theorem}

We have shown on the page that the listed generators \emph{are} integral; that there are no
further integral classes in $K\otimes\Q$ is the part we take from the literature. Granting
it, the numbers (ii)--(iii) can be checked by hand, which is a useful exercise in the lattice
techniques of \cref{ch:lattices}.

\begin{example}[The index $2^5$ and the discriminant $2^6$]\label{ex:kummer:index}
A vector of $K$ is $x=\frac12\sum_in_i\bar E_i$ with $n_i\in\Z$, and the condition is on the
parities: $(n_i\bmod2)\in\mathcal{C}$, where $\mathcal{C}\subseteq\mathbb{F}_2^{16}$ is the
linear span of the indicator functions of the hyperplanes. The indicator function of
$\{a\cdot\epsilon=c\}$ is the affine function $\epsilon\mapsto a\cdot\epsilon+c+1$ on
$\mathbb{F}_2^4$. The span is the space of all affine functions
$c_0+c_1\epsilon_1+\dots+c_4\epsilon_4$, of dimension $5$: $|\mathcal{C}|=2^5=32$ words,
namely $0$, the all-ones word, and the $30$ hyperplanes (weight $8$). The index of
$\bigoplus\Z\bar E_i$ in $K$ is $|\mathcal{C}|=2^5$. The index formula of \cref{ch:lattices},
$\operatorname{disc}(\text{sublattice})=\operatorname{disc}(\text{lattice})\cdot
\text{index}^2$, gives
\[
  \operatorname{disc}K=\frac{\operatorname{disc}\langle-2\rangle^{\oplus16}}{(2^5)^2}
  =\frac{2^{16}}{2^{10}}=2^6 .
\]
\emph{Integrality:} two hyperplanes meet in $0$, $4$ or $8$ points, so
$\frac12\sum_W\bar E_i\cdot\frac12\sum_{W'}\bar E_i=-\frac12|W\cap W'|\in\{0,-2,-4\}$.
\emph{Evenness:} $(\frac12\sum_W\bar E_i)^2=\frac14\cdot8\cdot(-2)=-4$. \emph{Roots:} a vector
with some half-integral $n_i/2$ has at least $8$ of them, hence square $\le-4$; the vectors
of square $-2$ are the $\pm\bar E_i$ only. \emph{Global check:} $H^2(X,\Z)$ is unimodular
and $K$, $K^\perp$ are primitive, so $|\operatorname{disc}K|=|\operatorname{disc}K^\perp|$
\source{Huy ll.~13512--13515}: indeed $2^6=|-2^6|$ by \eqref{eq:kummer:U2}. The index of
$K^\perp\oplus K$ in $H^2(X,\Z)$ is $\sqrt{2^6\cdot2^6}=2^6$, and that of the naive lattice
$U(2)^{\oplus3}\oplus\langle-2\rangle^{\oplus16}$ is $2^6\cdot2^5=2^{11}
=\sqrt{2^6\cdot2^{16}}$. (An explicit basis of $K$ was built and its Gram determinant $2^6$,
integrality and evenness were confirmed by a symbolic check; the code $\mathcal{C}$ was
enumerated by machine: $32$ words of weights $0,8,16$.) The remaining factor $2^6$ is made of
``mixed'' classes such as \eqref{eq:kummer:C}, which glue $K^\perp$ to $K$.
\end{example}

\begin{pitfallbox}[$22=6+16$ is not a lattice decomposition]
It is tempting to write $H^2(\operatorname{Km}(A),\Z)=U^{\oplus3}\oplus
\langle-2\rangle^{\oplus16}$. The right-hand side has discriminant $-2^{16}$ and is not
unimodular; with $U(2)$ in place of $U$ it is $-2^{22}$. The true lattice
$3U\oplus2E_8(-1)$ contains $U(2)^{\oplus3}\oplus\langle-2\rangle^{\oplus16}$ with index
$2^{11}$. In particular $\sum_i\bar E_i^2=-32$ is the trace of a Gram matrix of a sublattice in
one basis, not an invariant of $H^2(X,\Z)$.
\end{pitfallbox}

\subsection{Which K3 surfaces are Kummer?}

The Picard lattice of $X=\operatorname{Km}(A)$ contains the $\bar E_i$ and the images of the
algebraic classes of $A$, and the transcendental lattices are related by the rescaling
\eqref{eq:kummer:factor2}:
\begin{equation}\label{eq:kummer:rho}
  \rho(X)=\rho(A)+16,\index{Picard number!of a Kummer surface}\qquad T(X)\cong T(A)(2),\qquad 16\le\rho(X)\le20\ \ (\text{over }\C)
\end{equation}
\source{Huy §3.2.5, ll.~2084--2095; Asp ll.~1192--1194}. A Kummer surface need not be
algebraic, because a complex torus need not be \source{Asp ll.~1193--1194}; for a generic
torus $\rho(A)=0$ (standard). At the other extreme, the square torus of \cref{ex:kummer:square} is
$E_{\ii}\times E_{\ii}$ with $E_{\ii}$ the elliptic curve with complex multiplication by
$\ii$; it has $\rho(A)=4$ (the classes of the two factors, the diagonal and the graph of
multiplication by $\ii$; a standard fact, not checked against this book's library) and its
Kummer surface has the maximal Picard number $20$.

\begin{theorem}[Nikulin; \TL]\label{thm:kummer:nikulin}\index{Nikulin's theorem}
For a complex K3 surface $X$ the following are equivalent: (i) $X$ is a Kummer surface;
(ii) there is a primitive embedding $K\hookrightarrow NS(X)$. Both imply (iii): there is a
primitive embedding $T(X)\hookrightarrow U(2)^{\oplus3}$; and (iii) implies (i) if $X$ is
projective \source{Huy Thm.~14.3.17, ll.~13545--13552}. Moreover $X$ is Kummer if and only
if it contains $16$ disjoint smooth rational curves \source{Huy Rem.~14.3.19,
ll.~13595--13597}.
\end{theorem}

The proof combines the global Torelli theorem of \cref{ch:k3lattice} with the surjectivity
of the period map for tori, applied to $K^\perp\cong U(2)^{\oplus3}$
\source{Huy ll.~13578--13586}. As an application, the Fermat quartic
$x_0^4+x_1^4+x_2^4+x_3^4=0$ of \cref{ch:k3} has $T(X)\cong\Z(8)\oplus\Z(8)$, which embeds in
$U(2)^{\oplus3}$; hence it is a Kummer surface, although no torus is visible in its equation
\source{Huy Example~14.3.18, ll.~13588--13593}.

\begin{historybox}
The term ``orbifold'' is due to W.~Thurston; the objects first appeared in the mathematical
literature as ``V-manifolds'' \source{Asp ll.~1036--1039}. \Cref{thm:kummer:nikulin} was
first stated by Nikulin \source{Huy l.~13544}, and the global Torelli theorem was first
proved for Kummer surfaces \source{Huy ll.~13579--13581}.
\end{historybox}

%=====================================================================================
\section{Strings on the orbifold: twisted sectors}\label{sec:kummer:strings}

A point particle on $S_0$ sees a singular space: its wave equation needs boundary conditions
at the $16$ cone points that the geometry does not supply. The string does better. String
theory on $A/\iota$ is a perfectly regular conformal field theory, built from the free theory
on the torus by a three-step prescription \source{GPR §3.1, ll.~4118--4148}:
\index{orbifold!string theory on}
\begin{enumerate}[leftmargin=*,itemsep=2pt]
\item define the action of the group on all worldsheet fields, here $g\colon X^m\mapsto-X^m$
and, by worldsheet supersymmetry, $\psi^m\mapsto-\psi^m$ for the four torus directions
$m=1,\dots,4$;
\item keep only the invariant states, with the projector $P=\frac12(1+g)$;
\item add the strings that close only up to the group action, the \emph{twisted
sector}\index{twisted sector},
\begin{equation}\label{eq:kummer:twistedbc}
  X^m(\sigma+2\pi,\tau)=-X^m(\sigma,\tau),
\end{equation}
and project these onto invariant states as well.
\end{enumerate}
Step 3 has no analogue for particles. A string obeying \eqref{eq:kummer:twistedbc} is an open
curve on the torus whose end-points are exchanged by $\iota$, hence a \emph{closed} curve on
$S_0$; leaving it out would be as arbitrary as leaving out winding strings on a circle
(\cref{ch:circle}).

\subsection{The untwisted sector:\index{untwisted sector} \texorpdfstring{$16$}{16} moduli and no gauge fields}

The untwisted states are the states of the torus compactification of \cref{ch:narain} that
are even under $g$. A state with momentum and winding $(n_m,w^m)$ is sent to
$(-n_m,-w^m)$, so only the combinations $\ket{n,w}+\ket{-n,-w}$ (times an even oscillator
polynomial) and $\ket{n,w}-\ket{-n,-w}$ (times an odd one) survive. For the massless
NS--NS fields the rule is: count internal indices.
\begin{itemize}[leftmargin=*,itemsep=2pt]
\item The scalars $G_{mn}$ and $B_{mn}$, with vertex operators
$\partial X^m\bar\partial X^n$, have two internal indices and are even. All
$10+6=16$ of them survive. They parametrize the Narain moduli space
$O(4,4)/(O(4)\times O(4))$ of the torus, of dimension $4\cdot4=16$.
\item The $4+4$ Kaluza--Klein and winding gauge fields $G_{\mu m}$, $B_{\mu m}$ have one
internal index and are odd. They are projected out. This is the string's way of saying
$b_1(X)=0$: a K3 surface has no continuous isometries and no one-cycles for $B$ to wrap.
\end{itemize}
The T-duality group $\Odd{4}$ of the torus (\cref{ch:odd}) commutes with $g=-\mathbf 1$ and
acts on the untwisted sector; it will reappear as a subgroup of the duality group
$O(\Gamma_{4,20})$ of strings on K3 in \cref{ch:stringsk3}.

\subsection{The twisted sector: \texorpdfstring{$16\times4$}{16 x 4} moduli}

The boundary condition \eqref{eq:kummer:twistedbc} has three consequences.

\emph{(i) Localization.} Writing $X^m=x^m+\text{oscillators}$, the constant part must
satisfy $x\equiv-x$ modulo $\Lambda$: the centre of mass sits at one of the $16$ fixed points
\eqref{eq:kummer:fixed}. There is one copy of the twisted Hilbert space for each fixed point
\source{GPR ll.~4362--4369}, and a twisted string cannot carry momentum or winding on the
torus \source{GPR ll.~4379--4380}. Twisted strings are stuck at the singularities; they are
the degrees of freedom of the sixteen spheres of zero size.

\emph{(ii) Moding.} An antiperiodic field has half-integer modes,
$X^m\sim\sum_{r\in\Z+1/2}\alpha^m_r\,\ee^{-\ii r(\tau-\sigma)}/r+\dots$\,. By worldsheet
supersymmetry $\psi^m$ picks up the same extra sign: in the twisted NS sector the four
$\psi^m$ are \emph{integer} moded, in the twisted R sector half-integer moded.

\emph{(iii) Zero-point energy.}\index{zero-point energy} Use the mnemonic of \source{Pol ll.~165--196}: a periodic
boson contributes $-\frac1{24}$, an antiperiodic boson $+\frac1{48}$, and fermions
contribute with the opposite sign. In light-cone gauge there are $8$ transverse directions,
$4$ along the uncompactified space and $4$ along the torus. For one side (left- or
right-moving) of the NS sector
\begin{equation}\label{eq:kummer:E0}
  E_0^{\text{untw}}=8\Bigl(-\tfrac1{24}-\tfrac1{48}\Bigr)=-\tfrac12,\qquad
  E_0^{\text{tw}}=4\Bigl(-\tfrac1{24}-\tfrac1{48}\Bigr)+4\Bigl(\tfrac1{48}+\tfrac1{24}\Bigr)
  =-\tfrac14+\tfrac14=0 .
\end{equation}
This is the same arithmetic as for an open string with $\nu$ mixed Neumann--Dirichlet
directions in \cref{ch:dbranes}, $E_0=-\frac12+\frac\nu8$ \source{Pol eq.~(98),
ll.~2160--2180}, at $\nu=4$. In the untwisted sector the ground state is tachyonic and is
removed by the GSO projection, and the massless states need one oscillator $\psi_{-1/2}$. In
the twisted sector the \emph{ground state itself} is massless. It is degenerate, because the
four integer-moded $\psi^m$ have zero modes $\psi^m_0$ obeying a Clifford algebra: $2^{4/2}=4$
states, a spinor of the internal $SO(4)$, of which the GSO projection keeps one chirality,
$2$ states. Combining left- and right-movers,
\begin{equation}\label{eq:kummer:4scalars}
  2\times2=4\ \text{massless NS--NS scalars at each fixed point}\;=\;\mathbf 3\oplus\mathbf 1 .
\end{equation}
(The zero-point energies are from the cited source. The counting of the GSO-projected
zero-mode states and their survival under $P$ is the standard result for this orbifold; it
is not checked against a source in this book's library.) The decomposition
$\mathbf3\oplus\mathbf1$ refers to the $SU(2)$ rotating the three-plane $\Sigma$. The triplet
is the three metric parameters of the physics box in \cref{sec:kummer:local}: giving it an
expectation value blows up or deforms the $A_1$ point. The singlet is the flux of the
$B$-field through the sphere, $\int_{\bar E_i}B$.

\begin{example}[$80=16+64$, and two refinements]\label{ex:kummer:80}
The untwisted sector has $16$ moduli and the twisted sectors $16\times4=64$, together
\[
  16+64=80=4\cdot20=\dim\frac{O(4,20)}{O(4)\times O(20)} ,
\]
the dimension of the moduli space of conformal field theories on K3
\source{Asp ll.~1670 and 1948}. The count can be split according to geometric origin
\source{Asp ll.~1650--1670}:
\[
 \underbrace{10}_{\text{flat metrics on }T^4}+\underbrace{16\times3}_{\text{triplets}}=58
 \ \ \text{(Einstein metrics on K3)},\qquad
 \underbrace{6}_{B_{mn}}+\underbrace{16\times1}_{\text{singlets}}=22=b_2\ \ \text{($B$-field)}.
\]
The number $58$ is $3\cdot19+1$, the dimension of the space of positive three-planes in
$\R^{3,19}$ plus the volume (\cref{ch:k3lattice}); the number $22=6+16$ is
\eqref{eq:kummer:b2} once more. The massless spectrum of the singular orbifold thus knows the
Betti number of the smooth K3 surface. In M-theory language the $16$ twisted two-forms
come with $16$ of the $22$ gauge fields obtained by reducing the three-form on $H^2(X)$
\source{DRS ll.~511--519}.
\end{example}

\begin{pitfallbox}[The solvable orbifold is not the point of enhanced gauge symmetry]\index{enhanced gauge symmetry}\index{B-field@$B$-field!at orbifold points}
When all sixteen spheres shrink, M-theory on $T^4/\Z_2$ has the non-abelian gauge group
$SU(2)^{16}\times U(1)^6$, one $SU(2)$ per $A_1$ point \source{DRS ll.~515--519}, and for the
type IIA string an $A$--$D$--$E$ singularity likewise gives the corresponding gauge group
\source{Asp Prop.~2, ll.~2517--2529}. In the perturbative type IIA string, however, no new
massless vectors can appear \source{Asp ll.~2441--2443}, and the orbifold conformal field
theory is perfectly regular. There is no contradiction: a vanishing sphere is necessary but
not sufficient. The gauge symmetry is enhanced when, in addition, $B=0$ along the class of
the sphere, whereas the conformal field theory orbifold has $B=\frac12$ there. The
free-field orbifold and the point with gauge group $SU(2)^{16}$ are different points of the
$80$-dimensional moduli space, with the same singular metric \TL{}
\source{Asp ll.~2530--2547}. The singlet in \eqref{eq:kummer:4scalars} is the modulus that
interpolates between them.
\end{pitfallbox}

A classical observer sees sixteen cone points of infinite curvature. The string sees, at
each of them, four massless scalars whose expectation values are the size, the shape and the
$B$-flux of a sphere that is not yet there. Turning them on moves the theory into the smooth
Kummer surfaces, which are dense in the moduli space of K3 surfaces
\source{Huy ll.~6200--6202}: the solvable model is a point \emph{of} the K3 moduli space, not
a caricature of it (compare \cref{ch:dualscale} for the theme that at sizes of order
$\sqrt{\ap}$ ``singular'' and ``smooth'' are not geometric notions).

%=====================================================================================
\section{What the machine has checked}\label{sec:kummer:lean}

Three files of the repository contain declarations about the Kummer construction:
\begin{center}\small
\begin{tabular}{@{}ll@{}}
\lean{KummerBlowup.lean} & in \lean{StringTheoryFormalization/StringDynamics/}\\
\lean{KummerTadpole.lean} & in \lean{DualScaleM24Formalization/Moonshine/}\\
\lean{Problem9\_KummerModularity.lean} & in \lean{Lean5Corpus/Problems/}
\end{tabular}
\end{center}
All of them belong to the older part of the corpus that records \emph{arithmetic shadows}:
integers and one literal matrix.
None of them defines a torus, an involution, a blow-up, a cohomology group or a lattice
embedding, so nothing in the preceding sections is Tier~A as a statement about surfaces.
What is kernel-checked is the arithmetic that those sections reduce to. The theorems below depend
only on the standard axioms \lean{propext}, \lean{Classical.choice} and \lean{Quot.sound}.

\begin{leanbox}[The $16\times16$ Gram matrix (\lean{KummerBlowup.lean}, namespace
\lean{StringTheory.StringDynamics})]
\begin{leancode}
def kummerExceptionalDivisors : Fin 16 → String :=
  fun i => s!"E_{i.val}"

def kummerIntersectionForm (i j : Fin 16) : ℤ :=
  if i = j then -2 else 0

theorem kummer_lattice_contribution :
    (∑ i : Fin 16, kummerIntersectionForm i i) = -32 := by ...

theorem exceptional_self_intersection (i : Fin 16) :
    kummerIntersectionForm i i = -2 := by ...
\end{leancode}
\textbf{Reading.} \lean{kummerIntersectionForm} is the matrix $-2\cdot\mathbf 1_{16}$, the
Gram matrix of $\langle-2\rangle^{\oplus16}$ in the basis $\bar E_i$. The first definition
only produces the names \lean{"E\_0"},\dots,\lean{"E\_15"}. The theorem
\lean{exceptional\_self\_intersection} says that every diagonal entry of this matrix is $-2$;
\lean{kummer\_lattice\_contribution} says that its trace is $-32$. They are statements about
a matrix that was \emph{defined} to have these entries. The mathematical content of
\cref{prop:kummer:minus2}, that the curve produced by the blow-up has normal bundle
$\mathcal{O}(-2)$, is an input, not an output. Note also that the matrix describes
$\bigoplus\Z\bar E_i$, not the Kummer lattice $K$ of \cref{def:kummer:K}, which is larger by
index $2^5$; and by the pitfall box of \cref{sec:kummer:lattice} the number $-32$ is a trace in
one basis, not a lattice invariant. The companion \lean{kummer\_sublattice\_rank} in
\lean{MukaiLattice.lean} states that the index type \lean{Fin 16} has $16$ elements.

\textbf{Proof idea.} \lean{simp} unfolds the definition; on the diagonal the condition
\lean{i = i} is true and the entry is $-2$. For the trace, \lean{Finset.sum\_const} turns the
sum of a constant over \lean{Fin 16} into $16\cdot(-2)$.
\end{leanbox}

\begin{leanbox}[Counterparts in the other libraries]
\begin{leancode}
-- DualScaleM24Formalization/Moonshine/KummerTadpole.lean
--   (namespace SocrateAI.Moonshine)
def num_kummer_singularities : Nat := 16
def exceptional_self_intersection : Int := -2
theorem k3_euler_eq_24 : k3_euler = 24 := by rfl
theorem kummer_exceptional_intersection (i j : Fin 16) :
    (if i = j then exceptional_self_intersection else 0) =
    (if i = j then -2 else 0) := by rfl
theorem picard_rank_kummer_maximal : k3_b2 - transcendental_rank = 20 := by rfl
theorem hodge_h11_eq_20 : k3_b2 - 2 * hodge_h20 = 20 := by rfl

-- Lean5Corpus/Problems/Problem9_KummerModularity.lean
theorem k3_euler_characteristic_identity :
    k3_betti_0 + k3_betti_2 + k3_betti_4 = k3_euler_char := by rfl
theorem kummer_cycle_balance :
    kummer_exceptional_divisors + kummer_torus_cycles = 20 := by rfl
theorem singular_k3_picard_number_equals_twenty :
    picard_number singular_kummer_shioda_tate = 20 := by rfl

-- StringTheoryFoundation/ModularForms/FermatModularBridge.lean
theorem kummer_fixed_points_dim4 :
    (2 : Nat) ^ 4 = kummer_fixed_points_count := by rfl
\end{leancode}
\textbf{Reading.} Here \lean{k3\_euler} is defined as $1-0+22-0+1$ from the constants
\lean{k3\_b0},\dots,\lean{k3\_b4}, so \lean{k3\_euler\_eq\_24} is the last equality of
\eqref{eq:kummer:b2}; \lean{k3\_euler\_characteristic\_identity} is the same sum $1+22+1=24$
with natural numbers. \lean{kummer\_cycle\_balance} is $16+4=20$, the count
$h^{1,1}=4+16$ below \eqref{eq:kummer:b2}; \lean{hodge\_h11\_eq\_20} is $22-2\cdot1=20$;
\lean{kummer\_fixed\_points\_dim4} is $2^4=16$, the cardinality in \eqref{eq:kummer:fixed}.
\lean{kummer\_exceptional\_intersection} compares two \lean{if}-expressions that become
identical once the constant $-2$ is unfolded. \lean{HodgeNumbers.lean} has one more
statement of $20=16+4$, \lean{hodge11\_from\_kummer}. The second half of \lean{KummerTadpole.lean}
(\lean{rr\_tadpole\_cancellation}: $16\cdot4+4\cdot(-16)=0$) concerns seven-brane charges and
is discussed in \cref{ch:tadpole}.

\textbf{Two names that claim more than the statement.}
\lean{picard\_rank\_kummer\_maximal} is $22-2=20$. By \eqref{eq:kummer:rho} the Picard number
of a Kummer surface is $16+\rho(A)$ and takes every value from $16$ to $20$; the value $20$
belongs to special tori such as that of \cref{ex:kummer:square}, not to Kummer surfaces in
general. In \lean{singular\_kummer\_shioda\_tate} the Shioda--Tate formula\index{Shioda--Tate formula}
$\rho=2+\rk MW+\sum_v(m_v-1)$ is instantiated with the triple $(2,0,18)$. For the natural
fibration of $\operatorname{Km}(E_1\times E_2)$ of \cref{ex:kummer:24}(c) the four $I_0^*$
fibres contribute $4\cdot4=16$, so that $\rho=18+\rk MW$ \source{Huy Example~11.3.5,
ll.~9881--9885}, and $\rho=20$ is reached with the triple $(2,2,16)$. Lean checks that
$2+0+18=20$; it does not know which fibration, if any, has these data.

\textbf{Proof idea.} All of these are closed by \lean{rfl} (or \lean{decide}): both sides
reduce to the same numeral by evaluation.
\end{leanbox}

\begin{center}\small
\begin{tabular}{@{}p{0.47\textwidth}p{0.06\textwidth}p{0.41\textwidth}@{}}
\toprule
Claim & Tier & Where \\
\midrule
$2^4=16$; trace of $-2\cdot\mathbf1_{16}$ is $-32$; $1+22+1=24$; $16+4=20$; $22-2=20$
 & \TA & declarations quoted above\\
These integers are the fixed-point number, $\sum\bar E_i^2$, $\chi$, $h^{1,1}$ and
$\rho_{\max}$ of a Kummer surface & \TL & identification; never Tier A\\
$\operatorname{Km}(A)$ is a K3 surface; $\bar E_i^2=-2$; $b_2=6+16$; Nikulin's criterion
 & \TL & Asp §2.6; Huy ch.~1, §3.2.5, Thm.~14.3.17 \\
$K$ has index $2^5$ over $\bigoplus\Z\bar E_i$, $\operatorname{disc}K=2^6$,
 $K^\perp\cong U(2)^{\oplus3}$ & \TL & Huy Prop.~14.3.14; numbers re-derived in
 \cref{ex:kummer:index} (symbolic check)\\
Twisted sectors: $E_0=0$, $4$ scalars per fixed point, $16+64=80$ & \TL &
 Pol eq.~(98); GPR §3.1; Asp l.~1670; GSO count standard, unsourced here\\
$B=\frac12$ at the conformal field theory orbifold & \TL & Asp ll.~2530--2547\\
\bottomrule
\end{tabular}
\end{center}

\paragraph{What a faithful formalization would need.} The lattice part of this chapter is
within reach of the Gram-matrix approach of \cref{ch:lattices}: take the rational basis
matrix $P$ of \cref{ex:kummer:index} (eleven $\bar E_i$ and five half-sums) and let the kernel
check that $P^{\mathsf T}(-2\cdot\mathbf 1)P$ has integer entries, even diagonal and
determinant $2^6$; adding $U(2)^{\oplus3}$ and six glue vectors of the type
\eqref{eq:kummer:C}, one could check unimodularity of the full rank-$22$ Gram matrix. The
geometric part (that the blow-up of $A/\iota$ is a K3 surface) needs complex surfaces,
blow-ups and sheaf cohomology, which the libraries used in this book do not provide; see
\cref{ch:open}.

%=====================================================================================
\begin{summarybox}
\begin{itemize}[leftmargin=*,itemsep=2pt]
\item $x\mapsto-x$ on $T^4=\R^4/\Lambda$ has the $2^4=16$ fixed points
$\frac12\Lambda/\Lambda\cong\mathbb{F}_2^4$. The quotient is flat with holonomy
$\Z_2\subset SU(2)$, simply connected, with an invariant $\dd z_1\wedge\dd z_2$: a K3 orbifold.
\item Locally $\C^2/\Z_2=\{x_0x_1=x_2^2\}$. Its blow-up is the total space of
$\mathcal{O}_{\mathbb{P}^1}(-2)$, with coordinates $(a,t)$, $a'=t^2a$; the exceptional sphere
has $E^2=-2$ and $\dd\xi\wedge\dd\eta=\frac12\dd a\wedge\dd t$ has no zeros. The area
$J\cdot E$ is a modulus, and $J\cdot E\to0$ returns the orbifold.
\item $b_2=6+16=22$, $h^{1,1}=4+16=20$, signature $(3,3)+(0,16)=(3,19)$, and $\chi=24$ in three
ways: $1+22+1$; $\frac{0-16}2+32$; $4\times\chi(I_0^*)=4\times6$.
\item Integrally, $\pi_*H^2(A,\Z)\cong U(2)^{\oplus3}$ and the $\bar E_i$ span
$\langle-2\rangle^{\oplus16}$, but $H^2(X,\Z)$ is larger by index $2^{11}$. The Kummer lattice
$K\supset\bigoplus\Z\bar E_i$ (index $2^5$, discriminant $2^6$) is generated by half-sums over
affine hyperplanes of $\mathbb{F}_2^4$, which come from pairs of parallel invariant tori. A K3
surface is Kummer iff $K$ embeds primitively in $NS(X)$ (Nikulin).
\item The string on $T^4/\Z_2$ has $16$ untwisted moduli and, in each of $16$ twisted sectors
with zero-point energy $0$, four massless scalars $\mathbf3\oplus\mathbf1$ (blow-up and
$B$-flux): $16+64=80=\dim O(4,20)/(O(4)\times O(20))$. The solvable orbifold has
$B=\frac12$ through each vanishing sphere and no enhanced gauge symmetry.
\item Lean checks the arithmetic ($2^4=16$, trace $-32$, $1+22+1=24$, $16+4=20$); the
geometry and the Kummer lattice are Tier~L.
\end{itemize}
\end{summarybox}

\section*{Exercises}
\addcontentsline{toc}{section}{Exercises}

\begin{exercise}[$\star$ Fixed points in other dimensions]\label{exr:kummer:fixed}
Show that $x\mapsto-x$ on $T^d$ has $2^d$ fixed points. For $d=2$ verify the Gauss--Bonnet
theorem for the pillow. For $d=4$ explain why $\chi(T^4/\Z_2)=8$ is not $\frac12\chi(T^4)$.
\end{exercise}

\begin{exercise}[$\star$ The quadric cone]\label{exr:kummer:cone}
(a) Show that every polynomial in $(\xi,\eta)$ that is even under
$(\xi,\eta)\mapsto(-\xi,-\eta)$ is a polynomial in $x_0,x_1,x_2$ of
\eqref{eq:kummer:invariants}. (b) Show that $f=x_0x_1-x_2^2$ is singular only at the origin.
(c) Find the linear change of variables that brings $f$ to $u^2+v^2+w^2$, and show that the
deformed surface $u^2+v^2+w^2=\varepsilon$, $\varepsilon>0$, contains the real two-sphere of
radius $\sqrt\varepsilon$, which shrinks to the singular point as $\varepsilon\to0$.
\end{exercise}

\begin{exercise}[$\star\star$ The $I_0^*$ fibre]\label{exr:kummer:I0}
With $F=2C+\sum_{i\in P}\bar E_i$ as in \cref{lem:kummer:C}, verify $F^2=0$,
$F\cdot C=0$ and $F\cdot\bar E_i=0$ from the intersection numbers \eqref{eq:kummer:C}. Write
the $5\times5$ Gram matrix of $(C,\bar E_i)$ and show that it is negative semi-definite
with kernel spanned by $F$. Deduce the contribution $4$ of each fibre to the Shioda--Tate
formula.
\end{exercise}

\begin{exercise}[$\star\star$ Discriminants and indices]\label{exr:kummer:disc}
Compute the discriminants of $U^{\oplus3}\oplus\langle-2\rangle^{\oplus16}$,
$U(2)^{\oplus3}\oplus\langle-2\rangle^{\oplus16}$ and $U(2)^{\oplus3}\oplus K$. Using the
index formula of \cref{ch:lattices}, find the index of the latter two in the unimodular lattice
$H^2(X,\Z)$ and check $2^{11}=2^5\cdot2^6$. Why can the first lattice not be a finite-index
sublattice of $H^2(X,\Z)$ in which the $\bar E_i$ are the exceptional curves?
\end{exercise}

\begin{exercise}[$\star\star\star$ Planes are not enough]\label{exr:kummer:planes}
Let $P\subset\mathbb{F}_2^4$ be an affine plane (four points). (a) Show that
$v=\frac12\sum_{i\in P}\bar E_i$ has even square and integral pairing with every element of
$K$, so that integrality alone does not exclude $v$ from $H^2(X,\Z)$. (b) Show that
nevertheless $v\notin H^2(X,\Z)$ when $P$ is the plane of an invariant torus $T$: otherwise
\cref{lem:kummer:C} would give $\frac12\pi_*[T]\in H^2(X,\Z)$, contradicting the primitivity
of $\pi_*H^2(A,\Z)$ stated in \cref{thm:kummer:K}(i). (c) Conclude that $v\in K^*\setminus K$
and find its order in $A_K$.
\end{exercise}

\begin{exercise}[$\star\star$ Zero-point energies]\label{exr:kummer:E0}
Generalize \eqref{eq:kummer:E0} to a $\Z_2$ reflection of $\nu$ of the eight transverse
directions and recover $E_0=-\frac12+\frac\nu8$. For which $\nu$ is the twisted NS ground
state massless? Show that the twisted R sector has $E_0=0$ for every $\nu$.
\end{exercise}

\begin{exercise}[$\star$ Lean: the off-diagonal entries]\label{exr:kummer:leanoff}
State and prove in Lean that distinct exceptional curves are disjoint in the toy Gram matrix:
\begin{leancode}
theorem exceptional_disjoint (i j : Fin 16) (h : i ≠ j) :
    kummerIntersectionForm i j = 0 := by
  sorry  -- hint: simp [kummerIntersectionForm, h]
\end{leancode}
Then state and prove that the matrix is symmetric. (Hint: \lean{eq\_comm} inside the
\lean{if}.)
\end{exercise}

\begin{exercise}[$\star\star\star$ Lean: the code behind the Kummer lattice]
\label{exr:kummer:leancode}
Model $\mathbb{F}_2^4$ as \lean{Fin 4 → Bool} and an affine function by its five
coefficients. Prove by \lean{decide} that every non-constant affine function takes the value
\lean{true} exactly $8$ times, and that two of them which are neither equal nor
complementary are both \lean{true} on exactly $4$ points. Relate this to
\cref{ex:kummer:index}.
\end{exercise}

\section*{Further reading}
\addcontentsline{toc}{section}{Further reading}
\begin{itemize}[leftmargin=*,itemsep=1pt]
\item \source{Asp §2.6, ll.~1031--1230 and 1503--1530}: orbifolds, the blow-up of
$\C^2/\Z_2$ as used here, $(-2)$-curves, Theorem~4, resolution versus deformation;
\source{Asp §4.3, ll.~2440--2480 and 2500--2547}: enhanced gauge symmetry and the $B$-field
at orbifold points (\cref{ch:stringsk3}).
\item \source{Huy ch.~1, ll.~218--257}: the algebraic construction and
$\omega_X\cong\mathcal{O}_X$; \source{Huy §3.2.5, ll.~2015--2096}: cohomology, the factor $2$,
$\rho(X)=\rho(A)+16$; \source{Huy §14.3.3, ll.~13460--13600}: the Kummer lattice and Nikulin's
theorem; \source{Huy ch.~11, ll.~9500--9507, 9725--9735, 9881--9885}: the $I_0^*$ fibration.
\item \source{GPR §3.1, ll.~4118--4148, 4320--4332, 4362--4380}: the orbifold prescription and
twisted sectors (for a $\Z_3$ heterotic orbifold); \source{Pol ll.~165--196 and eq.~(98),
ll.~2160--2180}: zero-point energies.
\item \source{DRS §2, ll.~511--541}: $T^4/\Z_2$ in M-theory, the grouping $16+4+2$ of the
two-forms, and the type IIB orientifold used in \cref{ch:tadpole}.
\end{itemize}
